Fix \(n\) and \(P\in\mathcal P_{\mathrm G,n}^{\mathrm{JMS}}\).  The support restriction in \(\text{def:gaussian-class}\) gives
\[
 |Y|\le C_q
 \quad P\text{-almost surely}. \tag{1}
\]
Because \(Y\) is bounded, all conditional expectations below exist.  Conditional Jensen's inequality applied first to \(\sigma(X)\) and then to \(\sigma(X,T)\) yields
\[
 |q_0(X)|
 =|E_P[Y\mid X]|
 \le E_P[|Y|\mid X]
 \le C_q, \tag{2}
\]
and, on writing \(m(X,T)=E_P[Y\mid X,T]\),
\[
 |m(X,T)|
 \le E_P[|Y|\mid X,T]
 \le C_q
 \quad P\text{-almost surely}. \tag{3}
\]

By the residual definitions,
\[
 Y=q_0(X)+\theta_0(P)\eta+\xi,
 \qquad
 \eta=T-g_0(X).
\]
Both \(q_0(X)\) and \(\eta\) are \(\sigma(X,T)\)-measurable.  The conditional-mean restriction \(E_P[\xi\mid X,T]=0\) therefore implies
\[
 m(X,T)=q_0(X)+\theta_0(P)\eta
 \quad P\text{-almost surely}. \tag{4}
\]

Suppose for contradiction that \(\theta_0(P)\ne0\).  Under the Gaussian-treatment restriction,
\[
 \theta_0(P)\eta
 \sim N\!\left(0,\theta_0(P)^2\sigma^2\right).
\]
The boundary condition \(\sigma>0\) makes this Gaussian nondegenerate.  Its positive tail is unbounded, so the event
\[
 A=\{\theta_0(P)\eta>2C_q\}
\]
has strictly positive probability.  Intersecting \(A\) with the probability-one event on which (2)--(4) hold does not change its positive probability.  On that intersection,
\[
 m(X,T)
 =q_0(X)+\theta_0(P)\eta
 \ge -C_q+\theta_0(P)\eta
 >C_q,
\]
contradicting (3).  Hence \(\theta_0(P)=0\).  Since \(n\) and \(P\) were arbitrary, the conclusion holds throughout \(\mathcal P_{\mathrm G,n}^{\mathrm{JMS}}\) for every \(n\).

Now suppose \(\mathcal P_{\mathrm G,n}^{\mathrm{JMS}}\ne\varnothing\).  The statistic
\[
 \widetilde\theta(O_1,\ldots,O_n;\bar g_n,\bar q_n)\equiv0
\]
is an admissible data-dependent rule in \(\text{def:minimax-risks}\).  For every \(P\) in the class, the first part gives
\[
 (\widetilde\theta-\theta_0(P))^2=0
 \quad P\text{-almost surely},
\]
and therefore
\[
 \sup_{P\in\mathcal P_{\mathrm G,n}^{\mathrm{JMS}}}
 E_P[(\widetilde\theta-\theta_0(P))^2]=0. \tag{5}
\]
Likewise \(|\widetilde\theta-\theta_0(P)|=0\) almost surely, and the defining set for the generalized quantile contains \(t=0\); thus
\[
 Q_{P,1-\gamma}(|\widetilde\theta-\theta_0(P)|)=0
\]
for every class member and
\[
 \sup_{P\in\mathcal P_{\mathrm G,n}^{\mathrm{JMS}}}
 Q_{P,1-\gamma}(|\widetilde\theta-\theta_0(P)|)=0. \tag{6}
\]
Both minimax criteria are nonnegative by definition.  Equations (5) and (6), followed by the infimum over all statistics, prove
\[
 \mathfrak R_n^{\mathrm G}
 =\mathfrak M_{n,1-\gamma}^{\mathrm G}
 =0.
\]
Thus the zero statistic is minimax for both losses.

No overlap, positivity, rank, division, or inverse condition is used.  The finiteness of \(C_q\) supplies the support bound, \(\sigma>0\) supplies nondegenerate Gaussian support, and class nonemptiness is needed only for the displayed minimax statement to avoid dependence on a convention for suprema over the empty set.